%!TEX program = xelatex
\documentclass[a4paper]{article}
\usepackage{fontspec}
\setmainfont{Times New Roman}
\setsansfont{Arial}
\setmonofont{Consolas}[Scale=0.88]

\usepackage{tcolorbox}
\newtcolorbox{Qbox}[1]{colback=green!5!white,colframe=green!75!black,fonttitle=\bfseries,title=#1}

% Cấu hình hiển thị code Python
\usepackage{listings}
\usepackage{xcolor}

% Màu nền và màu sắc cho code
\definecolor{codebg}{rgb}{0.96, 0.98, 0.96}
\definecolor{keywordcolor}{rgb}{0.58, 0, 0.82}   % tím
\definecolor{stringcolor}{rgb}{0.0, 0.5, 0}      % xanh lá
\definecolor{commentcolor}{rgb}{0.0, 0.4, 0.4}   % xanh đậm
\definecolor{numbercolor}{rgb}{0.5, 0.5, 0.5}    % xám cho số dòng
\definecolor{frameborder}{rgb}{0.2, 0.2, 0.2}    % viền khung

% Style cho Python
\lstdefinestyle{pythonstyle}{
    language=Python,
    backgroundcolor=\color{codebg},
    basicstyle=\ttfamily\footnotesize,
    keywordstyle=\color{keywordcolor}\bfseries,
    stringstyle=\color{stringcolor},
    commentstyle=\color{commentcolor}\itshape,
    numberstyle=\tiny\color{numbercolor},
    numbers=left,
    numbersep=8pt,
    frame=single,
    rulecolor=\color{frameborder},
    framerule=0.6pt,
    framesep=6pt,
    tabsize=4,
    showspaces=false,
    showstringspaces=false,
    breaklines=true,
    breakatwhitespace=true,
    captionpos=b,
    morekeywords={as, assert, async, await, nonlocal, True, False, None, self, fit, transform, fit_transform, predict, predict_proba}
}

\lstset{style=pythonstyle}

\usepackage{graphicx,wrapfig,lipsum}
\usepackage{float}
\usepackage{longtable}

%\usepackage[utf8]{inputenc}
%\usepackage[francais]{babel}
\usepackage{a4wide,amssymb,epsfig,latexsym,multicol,array,hhline,fancyhdr}
\usepackage{indentfirst}
\usepackage{amsmath}
\usepackage{amsfonts}
\usepackage{empheq}
\usepackage{lastpage}
\usepackage[lined,boxed,commentsnumbered]{algorithm2e}
\usepackage{enumerate}
\usepackage{color}
\usepackage{graphicx}							% Standard graphics package
\usepackage{array}
\usepackage{tabularx, caption}
\captionsetup{font={small,it}, position=below}
\usepackage{multirow}
\usepackage{multicol}
\usepackage{rotating}
\usepackage{graphics}
\usepackage{geometry}
\geometry{
    left=2cm,      % lề trái
    right=2cm,     % lề phải
    % top=2cm,         % lề trên
    % bottom=2cm,      % lề dưới
    % bindingoffset=0.5cm, % dùng nếu in 2 mặt (thêm vào lề trái)
}
\usepackage{setspace}
\usepackage{epsfig}
\usepackage{tikz}
\usepackage{refcount}
\usetikzlibrary{arrows,snakes,backgrounds}
\usepackage{hyperref}
\hypersetup{urlcolor=blue,linkcolor=black,citecolor=black,colorlinks=true} 
\usepackage{subfigure}
\usepackage{fancybox}
\newcolumntype{C}[1]{>{\centering\arraybackslash}p{#1}}
\newcolumntype{M}[1]{>{\centering\arraybackslash}m{#1}}
%\usepackage{pstcol} 								% PSTricks with the standard color package

\newtheorem{theorem}{{\bf Định lý}}
\newtheorem{property}{{\bf Tính chất}}
\newtheorem{proposition}{{\bf Mệnh đề}}
\newtheorem{corollary}[proposition]{{\bf Hệ quả}}
\newtheorem{lemma}[proposition]{{\bf Bổ đề}}


%\usepackage{fancyhdr}
\setlength{\headheight}{40pt}
\pagestyle{fancy}
\fancyhead{} % clear all header fields
\fancyhead[L]{
 \begin{tabular}{rl}
    \begin{picture}(25,15)(0,0)
    \put(-5,-15){\includegraphics[width=17mm, height=12mm]{hcmut.png}}
    %\put(0,-8){\epsfig{width=10mm,figure=hcmut.eps}}
   \end{picture}&
	%\includegraphics[width=8mm, height=8mm]{hcmut.png} & %
	\begin{tabular}{l}
		\textbf{Trường Đại Học Bách Khoa Tp.Hồ Chí Minh}\\
		\textbf{Khoa Khoa Học và Kỹ Thuật Máy Tính}
	\end{tabular} 	
 \end{tabular}
}
\fancyhead[R]{
	\begin{tabular}{l}
		\tiny \bf \\
		\tiny \bf 
	\end{tabular}  }
\fancyfoot{} % clear all footer fields
\fancyfoot[L]{\fontsize{9pt}{11pt}\selectfont Bài tập lớn môn Học Máy - CO3117 - 2026}
\fancyfoot[R]{\fontsize{9pt}{11pt}\selectfont Trang {\thepage}/\pageref{LastPage}}
\renewcommand{\headrulewidth}{0.1pt}
\renewcommand{\footrulewidth}{0.1pt}
%%%
\setcounter{secnumdepth}{4}
\setcounter{tocdepth}{3}
\makeatletter
\newcounter {subsubsubsection}[subsubsection]
\renewcommand\thesubsubsubsection{\thesubsubsection .\@alph\c@subsubsubsection}
\newcommand\subsubsubsection{\@startsection{subsubsubsection}{4}{\z@}%
                                     {-3.25ex\@plus -1ex \@minus -.2ex}%
                                     {1.5ex \@plus .2ex}%
                                     {\normalfont\normalsize\bfseries}}
\newcommand*\l@subsubsubsection{\@dottedtocline{3}{10.0em}{4.1em}}
\newcommand*{\subsubsubsectionmark}[1]{}
\makeatother


\begin{document}
\onehalfspacing
\begin{titlepage}

\begin{tikzpicture}[remember picture, overlay]

% Khung ngoài
\draw[line width=6pt]
  ([xshift=2cm,yshift=-2cm]current page.north west)
  rectangle
  ([xshift=-2cm,yshift=2cm]current page.south east);

% Khung trong
\draw[line width=3pt]
  ([xshift=2.3cm,yshift=-2.3cm]current page.north west)
  rectangle
  ([xshift=-2.3cm,yshift=2.3cm]current page.south east);

\end{tikzpicture}

\vspace{-1.5cm}

\begin{center}
\centering
\Large
\textbf{ĐẠI HỌC QUỐC GIA THÀNH PHỐ HỒ CHÍ MINH \\
TRƯỜNG ĐẠI HỌC BÁCH KHOA \\
KHOA KHOA HỌC VÀ KỸ THUẬT MÁY TÍNH }

\end{center}

\vspace{0.2cm}

\begin{figure}[h!]
\begin{center}
\includegraphics[width=7cm]{hcmut.png}
\end{center}
\end{figure}
\vspace{-0.7cm}
\begin{center}
		{\LARGE \bf BÁO CÁO} \\
        {\LARGE \bf BÀI TẬP LỚN HỌC MÁY}\\
			\vspace*{1cm}
			\begin{center}
				\fontsize{18pt}{15pt} \bf{Ứng dụng SVM, Naive Bayes, Logistic Regression, XGBoost} \\
                \vspace*{-.7cm}
                \fontsize{18pt}{30pt} \bf{vào bài toán phân loại review quần áo phụ nữ} \\
			\end{center}

\vspace{1cm}

    \begin{tabular}{lll}
    \Large \bf GVHD: & \Large Th.S Võ Thanh Hùng & 
    \vspace{0.2cm}\\
    \Large \bf SVTH1:  &\Large Nguyễn Quang Bảo &\Large 2410276\\
    \Large \bf SVTH2:  &\Large Cao Quang Đôn &\Large 2310752 \\
    \Large \bf SVTH3:  &\Large Trần Nhật Đăng &\Large 2310737 \\
    \end{tabular}
    \end{center}
    \vfill
    \vspace{0.5cm}
    \begin{center}
    {\large Thành phố Hồ Chí Minh, tháng 11/2026}
\end{center}
\end{titlepage}


%\thispagestyle{empty}
\newpage

\clearpage
\setcounter{page}{1} % bắt đầu lại từ trang 1
\newpage
\tableofcontents
\newpage
\listoffigures

\listoftables
\newpage

\input{"Content/Phan Tich/phan_tich.tex"}

\newpage

\input{"Content/Giai Quyet Bai Toan/giai_quyet_bai_toan.tex"}

\newpage

\input{"Content/So sanh, Phan Tich/so_sanh_ket_luan.tex"}

\end{document}